% Source: https://artofproblemsolving.com/community/c3h2749439_proving_pie_using_complementary_probabilities
% Stable identity: aops-topic-2749439-q01
\textbf{(a)} Since $P$ is monic with roots $p(A),p(B),p(C)$,
\[P(x)=\bigl(x-p(A)\bigr)\bigl(x-p(B)\bigr)\bigl(x-p(C)\bigr),\qquad P(1)=\bigl(1-p(A)\bigr)\bigl(1-p(B)\bigr)\bigl(1-p(C)\bigr).\]
By independence, $P(1)$ is the probability that none of the three events occurs, so the probability that at least one occurs is $1-P(1)$.

(Expanding $1-P(1)$ with Vieta's formulas gives $e_1-e_2+e_3$, where $e_1=\sum p(A)$, $e_2=\sum p(A)p(B)$, $e_3=p(A)p(B)p(C)$. This is exactly inclusion--exclusion for independent events.)

\textbf{(b)} $P(1)=1-\dfrac{13}{12}+\dfrac38-\dfrac1{24}=\dfrac{24-26+9-1}{24}=\dfrac14$, so the probability is $\boxed{\dfrac34}$. (Indeed $P(x)=\left(x-\dfrac12\right)\left(x-\dfrac13\right)\left(x-\dfrac14\right)$ and $1-\dfrac12\cdot\dfrac23\cdot\dfrac34=\dfrac34$.)
